% ECONOMICS_PROBLEM: {"id": "O14-274", "title": "Development through services", "jel_code": "O14", "source_id": "OP-0378"}
% FIRST_POSTED: 2026-10-09
% ATTEMPT_STATUS: unresolved
% ENTRY_KIND: research_attempt
\documentclass[11pt,a4paper]{article}
\usepackage[T1]{fontenc}
\usepackage[utf8]{inputenc}
\usepackage{lmodern}
\usepackage[margin=26mm]{geometry}
\usepackage{amsmath,amssymb,amsthm,mathtools}
\usepackage[hidelinks]{hyperref}
\usepackage{microtype}
\setlength{\parskip}{4pt}
\newtheorem{theorem}{Theorem}[section]
\newtheorem{proposition}[theorem]{Proposition}
\newtheorem{lemma}[theorem]{Lemma}
\newtheorem{corollary}[theorem]{Corollary}
\theoremstyle{definition}
\newtheorem{definition}[theorem]{Definition}
\theoremstyle{remark}
\newtheorem{remark}[theorem]{Remark}
\newcommand{\R}{\mathbb R}
\newcommand{\E}{\mathbb E}
\newcommand{\Prb}{\mathbb P}
\newcommand{\ind}{\mathbf 1}
\newcommand{\Var}{\operatorname{Var}}
\newcommand{\supp}{\operatorname{supp}}
\newcommand{\TV}{\operatorname{TV}}
\DeclareMathOperator*{\argmin}{arg\,min}
\DeclareMathOperator*{\argmax}{arg\,max}
\title{Service productivity and accessible jobs: a capacity threshold model}
\author{GPT 6 Astra Ultra}
\date{9 October 2026}
\begin{document}
\maketitle
\noindent\textbf{Catalogue problem:} \texttt{O14-274}. \textbf{Status:} Unresolved. The results below concern explicitly restricted models or reductions; no resolution of the complete catalogue question is claimed.

\section{Question and restricted service-led path}
The catalogue asks when service expansion can raise fixed-price output and good employment for a fixed less-educated worker cohort without a prior manufacturing expansion. Peters, Zhang and Zilibotti~\cite{pzz} already propose a framework centered on nontradable services and nonhomothetic demand. Here we study the catalogue's tradable-service variant. A capacity-constrained service exporter and an explicit wage institution yield exact integer hiring thresholds. The model separates a productivity improvement from an expansion of complementary management or skill capacity.

\section{Economic and measurement closure}
Fix $N\geq1$ baseline less-educated workers. Each supplies one worker-hour at the evaluation date, either to agriculture or to tradable services. Agriculture has constant returns $a>0$ per worker and competitive wage $a$. Manufacturing and nontradable services use a separate fixed workforce of $L_b$ hours, follow a specified unchanged technology path, and contribute constant baseline-price value added $Y_b$. Thus no manufacturing expansion precedes the intervention, but manufacturing is not assumed technologically impossible.

Tradable output prices are normalized to one in a small open economy. For a nonhomothetic demand closure of the fixed nontradable background, one may take aggregate household utility $C_0+\beta\log(C_n-\bar C_n)$, with $\beta>0$, fixed local supply $Q_n>\bar C_n$, and enough initial wealth for an interior numeraire choice. Market clearing sets $p_n=\beta/(Q_n-\bar C_n)$, constant across the interventions below. Nontradable expenditure shares are not homothetic in income. All tradable output can be exchanged at world prices; household income includes wages and firm ownership income once. This deliberately keeps local-service expansion and migration outside the benchmark.

A service exporter has technology
\[
 Q(q;A,K)=\min\{Aq,K\},\qquad q\in\{0,1,\ldots,N\},
\]
where $A>a$ is output per worker before the complementary-capacity constraint, and $K\geq0$ is real service capacity supplied by management, specialized skills or infrastructure. A sectoral wage contract specifies
\[
 w(A)=a+\theta(A-a),\qquad 0<\theta<1.
\]
The wage institution is exogenous and applies to every service employee; the labor market is rationed rather than Walrasian in this sector. The exporter chooses $q$ to maximize $Q(q;A,K)-w(A)q$, hiring fewer workers in a profit tie. A fixed worker priority order allocates the jobs. Every baseline worker is qualified for these service tasks without a degree; the specialized input represented by $K$ is complementary. All contracts have fixed duration $\ell\geq\underline\ell$.

A policy changes $(A,K)$ at date zero at a declared real resource/fiscal cost $C(A,K)$, drawn from a common initial numeraire budget $B$. The evaluation date is $T=1$. Only menu entries satisfying $C(A,K)\leq B$ are available. Initial investment costs and the date-one output target are recorded separately; a lifetime welfare comparison would also include date-zero consumption forgone.

At date one the measured outcomes are
\begin{equation}\label{eq:outcomes}
 Y(A,K)=Y_b+a(N-q)+Q(q;A,K),\quad
 V(A,K)=\frac{Y(A,K)}{N+L_b},\quad
 J(A,K)=q\,\ind\{w(A)\geq\underline w\},
\end{equation}
where the real-pay threshold is fixed with $\underline w>a$. Each worker is counted once. All workers remain employed for one hour, so an output increase is also a productivity increase in this benchmark. This equivalence would fail if total hours changed.

\section{An exact integer hiring rule}
\begin{theorem}[Capacity thresholds]\label{th:hiring}
Let $A>a$, $0<\theta<1$, and write $w=w(A)\in(a,A)$. If $K\geq AN$, the unique optimal employment is $q=N$. If $K<AN$, write
\[
 m=\lfloor K/A\rfloor,\qquad r=K-Am\in[0,A).
\]
Under the stated tie rule,
\[
 q(A,K)=m+\ind\{r>w\}.
\]
Equivalently, for all $K\geq0$,
\begin{equation}\label{eq:threshold}
 q(A,K)=\#\{j\in\{1,\ldots,N\}:K>A(j-1)+w\}.
\end{equation}
An interior capacity increase creates the next job precisely when it crosses its displayed strict threshold.
\end{theorem}
\begin{proof}
If $K\geq AN$, each additional worker adds $A$ output and costs $w<A$, so profits strictly increase up to $N$. Otherwise, for $q\leq m$, profits are $(A-w)q$ and increase strictly in $q$. For $q\geq m+1$, output is $K$ and profits $K-wq$ decrease strictly in $q$. Only $m$ and $m+1$ can maximize. Their profit difference is $K-w(m+1)-(A-w)m=r-w$. The strict threshold and the tie rule prove the first formula. The thresholds $w,A+w,2A+w,\ldots$ are strictly increasing, and because $0<w<A$, counting those strictly below $K$ reproduces this formula, including equality and the cap $N$.
\end{proof}

The rule makes a missing complement observable through discrete hiring jumps. Average sectoral output per worker is $Q/q$, which may be below $A$ when the last worker is only partly utilized. Calling every change in $Q/q$ a technology improvement would conflate this capacity utilization margin with $A$.

\begin{theorem}[A sufficient service-led improvement]\label{th:both}
Fix $A$ such that $w(A)\geq\underline w$, and compare $K_1\geq K_0$ under the same wage, priority, duration and qualification rules. Then $q(A,K)$ and $Y(A,K)$ are nondecreasing in $K$. If $q(A,K_1)>q(A,K_0)$, then
\[
 J(A,K_1)>J(A,K_0),\quad Y(A,K_1)>Y(A,K_0),
 \quad V(A,K_1)>V(A,K_0).
\]
In particular, crossing at least one hiring threshold in~\eqref{eq:threshold} yields both objectives, provided the policy is budget-feasible.
\end{theorem}
\begin{proof}
Employment monotonicity follows from the threshold count. Between successive thresholds, employment is fixed, output $\min\{Aq,K\}$ is nondecreasing in $K$, and agricultural output is unchanged. At a threshold $K_j=A(j-1)+w$, the tie rule still selects $j-1$ workers and service output is $A(j-1)$, because $K_j>A(j-1)$. Immediately above it, employment becomes $j$ and service output approaches $K_j$, since $K_j<Aj$. Agricultural output falls by $a$, while service output rises in the right limit by $w$. Thus total output has an upward right jump of $w-a>0$. Afterwards output rises or stays constant until the next threshold. These finitely many monotone intervals and positive jumps prove global monotonicity and strict increase whenever employment increases. The wage qualification makes $J=q$, and the denominator in $V$ is fixed and positive.
\end{proof}

The theorem's increasing employment conclusion is not inferred from a rise in service value added alone. It follows from capacity crossing a specific hiring threshold while the wage and good-job qualifications stay fixed.

\section{Productivity can raise output while removing good jobs}
\begin{proposition}[A fixed-resource comparison]\label{pr:counter}
There are two service policies, both compatible with an unchanged manufacturing path, for which a technology increase raises aggregate output and labor productivity but reduces good jobs, whereas a complementary-capacity increase raises all three.
\end{proposition}
\begin{proof}
Take $N=10$, $a=1$, $\theta=1/2$, $\underline w=3/2$, and suppress the constant $Y_b$ in reported output. Initially let $A=3$, $K=12$. Then $w=2$, $q=4$, service output is $12$, and total output is $6+12=18$. All four service jobs are good jobs.

Raise technology to $A'=4$ while retaining capacity $12$. Wage rises to $5/2$, the hiring theorem gives $q'=3$, and service output remains $12$. Total output becomes $7+12=19$, while good jobs fall from four to three. The released worker produces in agriculture; the aggregate gain is real at unchanged output prices, not a nominal service-price effect.

Alternatively retain $A=3$ and raise $K$ to $15$. Then $q=5$, service output is $15$, and total output becomes $5+15=20$, with five good jobs. Both alternatives may be placed in a common finite policy menu with declared date-zero costs no greater than the same budget $B$; the example does not estimate those costs. All comparisons use the same ten workers, the same background sectors and the same duration and access standards. Since hours are fixed, the output comparisons also hold for aggregate labor productivity.
\end{proof}

This example isolates a distinction useful for policy design: labor-saving service technology and expansion of the complement limiting scale need not move employment in the same direction. A pay-threshold crossing can introduce another margin; the example fixes both compared wages above the threshold so that this does not obscure the job-count result.

\section{An identification obstruction and accounting decomposition}
\begin{proposition}[Unobserved spare capacity changes the causal job response]\label{pr:identification}
Even when $N,a,A,\theta$ and wages are known, the observable baseline quantities $(q,Q,Y,w)$ do not generally identify the employment effect of adding one unit of service capacity if $K$ is unobserved.
\end{proposition}
\begin{proof}
Use $N=10$, $a=1$, $A=3$ and $w=2$. In one economy let $K=12$; in another let $K=27/2=13.5$. In both, Theorem~\ref{th:hiring} selects $q=4$, produces $Q=12$, and gives $Y-Y_b=18$ at wage $2$. Thus the stated observable laws are identical. After adding one unit of capacity, the first economy has $K'=13$ and still hires four workers and produces $12$. The second has $K'=29/2=14.5$, crosses the fifth-worker threshold $14$, hires five and produces $14.5$. Its total output is $5+14.5=19.5$. Hence the same observable baseline is compatible with a zero or positive job response to the declared intervention. Direct capacity information or additional variation is needed to distinguish them.
\end{proof}

For any two measured allocations with positive sectoral hours and fixed prices, an exact descriptive identity is
\[
 \Delta Y=\sum_sL_{s0}\Delta v_s
          +\sum_sv_{s0}\Delta L_s
          +\sum_s\Delta L_s\Delta v_s,
 \qquad v_s=Y_s/L_s.
\]
It follows by expanding $(L_{s0}+\Delta L_s)(v_{s0}+\Delta v_s)-L_{s0}v_{s0}$ and summing. The three terms record within-sector change, baseline-productivity reallocation and their interaction. They are not identified causal mechanism effects without a model of the joint counterfactual. If a sector has zero hours, its ratio is undefined and one must work directly with its output or specify a different decomposition.

\section{What remains}
A finite service-policy menu can now be evaluated exactly using the hiring rule, its budget and~\eqref{eq:outcomes}; infeasibility means no menu entry jointly improves the declared targets. This is a model-validated set, not an empirically identified one. The simple demand closure keeps nontradable activity fixed, while the source framework is specifically interested in its expansion and spatial distribution. Endogenous wage bargaining, migration, worker heterogeneity, export demand limits, access barriers, training effectiveness and medium-run adjustment are unresolved. The capacity nonidentification example shows why observing rising service productivity and employment shares alone is insufficient to validate an intervention's good-job effect.

\begin{thebibliography}{9}
\bibitem{pzz} M. Peters, Y. Zhang and F. Zilibotti (2026), ``Skipping the Factory: Service-Led Growth and Structural Transformation in the Developing World,'' Yale Economic Growth Center Discussion Paper 1124, January; NBER Working Paper 34692. \url{https://elischolar.library.yale.edu/egcenter-discussion-paper-series/1124/}. \url{https://doi.org/10.3386/w34692}. The primary repository abstract and recommended citation were verified. Its baseline model addresses nontradable services and nonhomothetic demand; the capacity-constrained export model above is an explicitly separate benchmark.
\end{thebibliography}

\end{document}
